% ===============================================================
%  جلسه ششم
% ===============================================================
\chapter{متغیر تصادفی و متغیرهای تصادفی گسسته معروف}

\section{متغیر تصادفی \lr{(Random Variable)} ــ مقدمه}

پیش‌تر دیدیم که مدل‌سازی تصادفی یک واقعه، معادل با تعریف یک فضای احتمال به صورت زیر است:
\[
(\Omega,\ \F,\ P).
\]
بیایید هر سه جزء را یک بار دیگر، این بار با نگاه «مهندسی»، مرور کنیم:

\begin{itemize}
\item \textbf{فضای نمونه $\Omega$:} مجموعه‌ای شامل تمامی \textbf{پیشامدهای} (\lr{outcome}) ممکن است. هر عضو این فضا مشخص‌کننده یک وضعیت خاص از سیستم مدل‌سازی شده است:
\[\omega=\big(\omega^{(1)},\omega^{(2)},\ldots\big)\in\Omega .\]
این فضا می‌تواند ابعاد بسیار بالایی داشته باشد، که به معنی زیاد بودن پارامترهای آزاد در مدل‌سازی است. $\Omega$ حتی می‌تواند \textbf{فضایی بی‌نهایت‌بعدی} باشد.
\item \textbf{فضای رویدادها $\F$:} مجموعه‌ای شامل تمام یا برخی زیرمجموعه‌های $\Omega$ است. وقوع هر پیشامد در $\Omega$، می‌تواند باعث رخ دادن هیچ، یک یا تعدادی (گاهاً بی‌نهایت) رویداد شود. لازم است تا $\F$ یک \lr{$\sigma$-Field} بر روی $\Omega$ باشد.
\item \textbf{قانون یا اندازه احتمال $P$:} تابعی از فضای رویدادها به بازه بسته اعداد حقیقی بین صفر تا یک است، یعنی
\[P:\F\to[0,1],\]
و لازم است تا $P$ قوانین کولموگروف \lr{(Kolmogorov)} را رعایت کند.
\end{itemize}

\subsection{یک مثال: مدل‌سازی آب‌وهوا}
فرض کنید قصد داریم وضعیت آب‌وهوای یک ناحیه از کره زمین را به صورت تصادفی مدل‌سازی کنیم:
\begin{itemize}
\item فضای نمونه در این حالت، بسیار پیچیده خواهد بود. لازم است تا محل و تکانه تمامی ذرات و مولکول‌های معلق در هوا را به عنوان پارامتر آزاد (بخشی از ابعاد هر نمونه $\omega\in\Omega$) در نظر بگیریم\ldots
\item همچنین، تمامی پارامترهایی که می‌توانند مستقلاً و یا در کنار یکدیگر بر مفهومی که ما از آب‌وهوا در ذهن داریم تأثیرگذار باشند.
\item رویدادهایی مانند «بارانی» (\lr{Rainy})، «آفتابی» (\lr{Sunny}) و «طوفانی» (\lr{Windy}) اعضای $\F$ هستند؛ هر کدام مجموعه بسیار بزرگی از وضعیت‌های میکروسکوپی $\omega$ است.
\end{itemize}

در عمل، \textbf{امکان سنجش فضای $\Omega$ اصلی وجود ندارد}. لذا، از اندازه‌گیری محدود و در دسترس استفاده خواهد شد:
\begin{itemize}
\item \textbf{دماسنج:} نگاشت فضای $\Omega$ به یک ناحیه از فضای $\R$ از طریق یک حسگر (\lr{sensor}) دمایی.
\item \textbf{رطوبت‌سنج:} اندازه‌گیری نادقیق تجمع مولکول‌های آب از طریق حسگر رطوبتی.
\item و \ldots
\end{itemize}
هر یک از این حسگرها یک تابع $X_1,X_2,\ldots$ از $\Omega$ به اعداد است.

\begin{center}
\begin{tikzpicture}[>=Stealth]
\node[draw=red!70!black,very thick,minimum width=3.2cm,minimum height=2cm,align=center] (O) at (0,0) {\large $\Omega$\\[2pt] \small (وضعیت کامل جو)};
\node[draw=blue!70!black,very thick,minimum width=3.2cm,minimum height=2cm,align=center] (R) at (7,0) {\large $\R$\\[2pt] \small (عدد خوانده‌شده)};
\draw[->,very thick] (O) -- node[above]{$X_1,X_2,\ldots$} node[below,font=\small]{حسگرها} (R);
\end{tikzpicture}
\end{center}

\subsection{اندازه‌گیری در فضای نمونه}
دستیابی به تمامی مقادیر لازم برای تشکیل فضای $\Omega$ و مشاهده یک پیشامد دلخواه $\omega\in\Omega$، بعضاً بسیار سخت و یا حتی غیرممکن است. لذا، به جای مشاهده (یا معادلاً اندازه‌گیری) خود $\omega$، تابعی از آن را مانند $X(\omega)$ اندازه می‌گیریم.

\begin{definition}[متغیر تصادفی]
یک \textbf{متغیر تصادفی} تابعی از فضای نمونه به یک مجموعه عددی است:
\[X:\Omega\to R .\]
در اینجا، برد تابع $X$ (یعنی $R$) یک \textbf{مجموعه عددی} دلخواه است. $R$ می‌تواند مجموعه اعداد حقیقی $\R$، اعداد طبیعی $\N$ یا مجموعه محدود $\{0,1\}$ باشد.
\end{definition}

\begin{note}
\notetitle نام «متغیر تصادفی» کمی گمراه‌کننده است: $X$ در واقع یک \textbf{تابع} (کاملاً معین) است و تصادف از انتخاب $\omega$ می‌آید. احتمال رویدادهایی مانند $\{X\le 3\}$ به صورت $P(\{\omega\in\Omega:\ X(\omega)\le3\})$ تعریف می‌شود؛ یعنی هر سؤال درباره $X$ در نهایت یک سؤال درباره یک رویداد در $\F$ است. (از نظر ریاضی لازم است این مجموعه‌ها عضو $\F$ باشند؛ به این شرط «اندازه‌پذیری» می‌گویند.)
\end{note}

\begin{example}[تعداد شیرها در دو پرتاب]
دو سکه سالم پرتاب می‌شوند: $\Omega=\{HH,HT,TH,TT\}$. متغیر تصادفی $X$ = «تعداد شیرها» را تعریف کنید. آنگاه
$X(HH)=2$، $X(HT)=X(TH)=1$ و $X(TT)=0$. پس
\[
P(X=0)=\tfrac14,\qquad P(X=1)=P(\{HT,TH\})=\tfrac12,\qquad P(X=2)=\tfrac14 .
\]
این مثال نشان می‌دهد که $X$ اطلاعات را «فشرده» می‌کند: دو پیشامد متفاوت $HT$ و $TH$ به یک مقدار نگاشته می‌شوند.
\end{example}

\subsection{متغیر تصادفی به مثابه یک «جعبه سیاه»}
برای فهم بهتر، می‌توانید یک متغیر تصادفی مانند $X$ را به مثابه یک \textbf{جعبه سیاه} در نظر بگیرید که نمی‌دانید در داخل آن چه خبر است\ldots
\begin{itemize}
\item هر بار که دکمه روی جعبه سیاه را فشار می‌دهید، یک آزمایش تصادفی درون آن انجام می‌گیرد (یک مقدار $\omega\in\Omega$ به صورت تصادفی مقدار می‌گیرد)، و یک مقدار تصادفی به $X$ می‌دهد.
\item به این مقدار یک \textbf{نمونه} از $X$ گفته می‌شود. در حالت کلی، هر بار که این کار را می‌کنید مقادیر مختلفی به دست خواهید آورد (مثلاً $1.2,\ 3.45,\ -2.1,\ 5,\ -0.4,\ldots$). در واقع $X$ یک مقدار مشخص ندارد، و در هر آزمایش یک مقدار خاص خواهد گرفت.
\item اما «\textbf{توزیع}» یا \lr{distribution} مقادیر آن مشخص است: اگر تعداد زیادی نمونه بگیریم و هیستوگرام آن‌ها را رسم کنیم، شکل هیستوگرام به یک شکل ثابت (توزیع $X$) میل می‌کند.
\end{itemize}

\section{متغیرهای تصادفی گسسته \lr{(Discrete Random Variables)}}

متغیرهای تصادفی می‌توانند پیوسته و یا گسسته باشند. \textbf{متغیرهای تصادفی گسسته} می‌توانند مقادیری از یک مجموعه محدود و یا حداکثر شمارا را داشته باشند.

\begin{definition}[تابع جرم احتمال \lr{(PMF)}]
برای متغیر تصادفی گسسته $X$ با مقادیر ممکن $x_1,x_2,\ldots$، تابع جرم احتمال عبارت است از
\[
P_X(x)=P(X=x),
\]
که در دو شرط صدق می‌کند: $P_X(x)\ge0$ و $\sum_i P_X(x_i)=1$.
\end{definition}

در صورتی که یک متغیر تصادفی گسسته باشد، \textbf{به طور کامل توسط \lr{PMF} تعیین خواهد شد}؛ زیرا برای هر مجموعه $A$ داریم
\[P(X\in A)=\sum_{x\in A}P_X(x).\]

\begin{example}[مجموع دو تاس چهاروجهی]
دو تاس چهاروجهی سالم (با وجوه ۱ تا ۴) را پرتاب می‌کنیم و $X$ را مجموع آن‌ها می‌گیریم. \lr{PMF} آن را بیابید.
\begin{solution}
۱۶ حالت هم‌احتمال داریم. تعداد حالت‌هایی که مجموع برابر $x$ می‌شود برای $x=2,\ldots,8$ به ترتیب $1,2,3,4,3,2,1$ است (مثلاً برای $x=5$: $(1,4),(2,3),(3,2),(4,1)$). پس:
\begin{center}
\renewcommand{\arraystretch}{1.3}
\begin{tabular}{c|ccccccc}
$x$ & $2$ & $3$ & $4$ & $5$ & $6$ & $7$ & $8$\\\hline
$P_X(x)$ & $\frac1{16}$ & $\frac2{16}$ & $\frac3{16}$ & $\frac4{16}$ & $\frac3{16}$ & $\frac2{16}$ & $\frac1{16}$
\end{tabular}
\end{center}
که جمع آن‌ها $\frac{16}{16}=1$ است.
\end{solution}
\end{example}

\begin{center}
\begin{tikzpicture}
\begin{axis}[width=6.5cm,height=4.5cm,ymin=0,ymax=0.6,xmin=-0.5,xmax=2.5,xtick={0,1,2},ytick={0.25,0.5},xlabel={$x$},title={\small $X$=تعداد شیر در دو پرتاب}]
\addplot+[ycomb,very thick,mark=*,teal] coordinates {(0,0.25) (1,0.5) (2,0.25)};
\end{axis}
\end{tikzpicture}
\hspace{5mm}
\begin{tikzpicture}
\begin{axis}[width=7.5cm,height=4.5cm,ymin=0,ymax=0.3,xmin=0.5,xmax=8.5,xtick={1,...,8},ytick={0.0625,0.125,0.1875,0.25},yticklabels={$\frac1{16}$,$\frac2{16}$,$\frac3{16}$,$\frac4{16}$},xlabel={$x$},title={\small مجموع دو تاس چهاروجهی}]
\addplot+[ycomb,very thick,mark=*,black] coordinates {(2,0.0625) (3,0.125) (4,0.1875) (5,0.25) (6,0.1875) (7,0.125) (8,0.0625)};
\end{axis}
\end{tikzpicture}
\end{center}

\section{متغیرهای تصادفی گسسته معروف}

\subsection{توزیع برنولی \lr{(Bernoulli)}}
ساده‌ترین متغیر تصادفی، نتیجه یک آزمون با دو خروجی «موفقیت» (۱) و «شکست» (۰) است. اگر احتمال موفقیت $p$ باشد، $X\sim\mathrm{Bernoulli}(p)$:
\[
P_X(1)=p,\qquad P_X(0)=1-p .
\]

\subsection{توزیع دوجمله‌ای \lr{(Binomial Distribution)}}
توزیع دوجمله‌ای مربوط به \textbf{تعداد موفقیت‌ها}ی تکرار $n$باره (مستقل) یک آزمون برنولی است.

\begin{theorem}[\lr{PMF} توزیع دوجمله‌ای]
برای $X\sim\mathrm{Binomial}(n,p)$ و $X=0,1,2,\ldots,n$:
\[
\boxed{\ P_X(k)=\binom nk p^k(1-p)^{n-k}\ }
\]
\end{theorem}
\begin{proof}
هر رشته مشخص از $n$ نتیجه که دقیقاً $k$ موفقیت دارد، به دلیل استقلال احتمال $p^k(1-p)^{n-k}$ دارد (ضرب احتمال‌ها). تعداد چنین رشته‌هایی $\binom nk$ است (انتخاب جای موفقیت‌ها) و این رشته‌ها ناسازگارند؛ پس با جمع کردن، حکم به دست می‌آید. جمع کل احتمال‌ها طبق قضیه دوجمله‌ای $(p+(1-p))^n=1$ است ــ و نام «دوجمله‌ای» از همین‌جاست.
\end{proof}

\begin{note}
\notetitle آزمایش سکه سالم در جلسه سوم حالت خاص $p=\frac12$ است: $\binom nk2^{-n}$.
\end{note}

\begin{center}
\begin{tikzpicture}
\begin{axis}[width=12cm,height=5cm,ymin=0,xmin=-0.5,xmax=20.5,xlabel={$k$},ylabel={$P_X(k)$},legend style={font=\small}]
\addplot+[ycomb,thick,mark=*,blue] table[x=k,y=a]{data/binom_p.dat};
\addlegendentry{\lr{$n=20,\ p=0.5$}}
\addplot+[ycomb,thick,mark=square*,red,xshift=2pt] table[x=k,y=b]{data/binom_p.dat};
\addlegendentry{\lr{$n=20,\ p=0.2$}}
\end{axis}
\end{tikzpicture}
\end{center}

\begin{example}[آزمون چندگزینه‌ای]
یک آزمون ۱۰ سؤال چهارگزینه‌ای دارد و دانشجویی همه سؤال‌ها را تصادفی جواب می‌دهد. احتمال اینکه دقیقاً ۵ سؤال را درست جواب دهد، و احتمال اینکه دست‌کم ۶ سؤال را درست جواب دهد چقدر است؟
\begin{solution}
تعداد پاسخ‌های درست $X\sim\mathrm{Binomial}(10,\frac14)$ است. پس
\[
P(X=5)=\binom{10}{5}\Big(\frac14\Big)^5\Big(\frac34\Big)^5=252\times\frac{243}{4^{10}}\approx0.0584,
\]
\[
P(X\ge6)=\sum_{k=6}^{10}\binom{10}{k}\Big(\frac14\Big)^k\Big(\frac34\Big)^{10-k}\approx0.0197 .
\]
یعنی شانس قبولی (با حد نصاب ۶) با حدس تصادفی کمتر از ۲٪ است.
\end{solution}
\end{example}

\subsection{توزیع چندجمله‌ای \lr{(Multinomial Distribution)}}
توزیع چندجمله‌ای، تعمیم توزیع دوجمله‌ای برای تکرار $n$باره آزمون‌هایی است که \textbf{بیش از دو خروجی} ممکن داشته باشند. فرض کنید آزمونی با $m$ خروجی ممکن، هر یک با احتمال‌های $p_1,p_2,\ldots,p_m$ (با جمع ۱)، را $n$ بار به صورت مستقل انجام می‌دهیم. هدف، به دست آوردن احتمال وقوع $k_i$ بار از پیشامد $i$ام است ($n=k_1+\cdots+k_m$).

\begin{theorem}[\lr{PMF} توزیع چندجمله‌ای]
برای $(k_1,\ldots,k_m)\sim\mathrm{Multinomial}(n,p_1,\ldots,p_m)$:
\[
\boxed{\ P_X(k_1,\ldots,k_m)=\binom{n}{k_1,\ldots,k_m}\prod_{i=1}^{m}p_i^{k_i}\ },\qquad
\binom{n}{k_1,\ldots,k_m}=\frac{n!}{k_1!\,k_2!\cdots k_m!}.
\]
\end{theorem}
\begin{proof}
هر رشته مشخص با $k_1$ بار خروجی ۱، \ldots\ و $k_m$ بار خروجی $m$ احتمال $\prod p_i^{k_i}$ دارد. تعداد این رشته‌ها برابر تعداد جایگشت‌های $n$ شیء است که $k_1$ تای آن‌ها یکسان (از نوع ۱)، $k_2$ تا یکسان و \ldots\ باشند؛ یعنی ضریب چندجمله‌ای $\frac{n!}{k_1!\cdots k_m!}$. (برای $m=2$ همان توزیع دوجمله‌ای به دست می‌آید.)
\end{proof}

\begin{example}[تاس و توزیع چندجمله‌ای]
(الف) تاس سالمی را ۶ بار می‌اندازیم. احتمال اینکه هر عدد دقیقاً یک بار ظاهر شود چقدر است؟
(ب) در یک انتخابات، ۵۰٪ به نامزد $A$، ۳۰٪ به $B$ و ۲۰٪ به $C$ رأی می‌دهند. اگر ۱۰ نفر را تصادفی انتخاب کنیم، احتمال اینکه ۵ نفر به $A$، ۳ نفر به $B$ و ۲ نفر به $C$ رأی داده باشند چقدر است؟
\begin{solution}
(الف) $n=6$، $m=6$، همه $k_i=1$ و $p_i=\frac16$:
$\ \frac{6!}{1!\cdots1!}\big(\frac16\big)^6=\frac{720}{46656}\approx0.0154$.

(ب) $\ \frac{10!}{5!\,3!\,2!}(0.5)^5(0.3)^3(0.2)^2=2520\times0.03125\times0.027\times0.04\approx0.0851$.
\end{solution}
\end{example}

\subsection{توزیع پواسون \lr{(Poisson Distribution)}}
متغیر تصادفی پواسون کاربرد فراوانی در مدل‌سازی \textbf{تعداد پیشامدهای} یک ماهیت تصادفی بی‌حافظه در طول زمان دارد (مثلاً تعداد تماس‌های ورودی به یک مرکز تلفن در یک ساعت، تعداد بسته‌های رسیده به یک روتر در یک ثانیه، تعداد ذرات تابیده‌شده از یک ماده رادیواکتیو). این توزیع بر روی اعداد حسابی $k=0,1,2,\ldots$ تعریف شده و تنها دارای یک پارامتر آزاد $\lambda\in\R_{\ge0}$ (نرخ) است.

\begin{definition}[توزیع پواسون]
برای $X\sim\mathrm{Pois}(\lambda)$:
\[
\boxed{\ P_X(k)=e^{-\lambda}\frac{\lambda^k}{k!},\qquad k=0,1,\ldots\ }
\]
\end{definition}
بررسی نرمال بودن: با بسط تیلور $e^\lambda=\sum_{k\ge0}\frac{\lambda^k}{k!}$ داریم $\sum_kP_X(k)=e^{-\lambda}e^{\lambda}=1$.

\begin{center}
\begin{tikzpicture}
\begin{axis}[width=12cm,height=5.5cm,ymin=0,xmin=-0.5,xmax=20.5,xlabel={$k$},ylabel={$P(X=k)$},legend style={font=\small}]
\addplot[orange,thick,mark=*,mark options={fill=orange}] table[x=k,y=p]{data/pois1.dat};
\addlegendentry{$\lambda=1$}
\addplot[violet,thick,mark=*,mark options={fill=violet}] table[x=k,y=p]{data/pois4.dat};
\addlegendentry{$\lambda=4$}
\addplot[cyan!70!black,thick,mark=*,mark options={fill=cyan}] table[x=k,y=p]{data/pois10.dat};
\addlegendentry{$\lambda=10$}
\end{axis}
\end{tikzpicture}
\end{center}

\subsection{خواص توزیع پواسون: حد توزیع دوجمله‌ای}
توزیع پواسون نتیجه یکی از حالت‌های توزیع دوجمله‌ای (\lr{Binomial}) است.

\begin{theorem}[تقریب پواسون]
در توزیع دوجمله‌ای $\mathrm{Binomial}(n,p)$، فرض کنید $n\to\infty$ و $p\to0$ به طوری که حاصل‌ضرب آن‌ها عددی ثابت بماند، یعنی $np=\lambda$. آنگاه برای هر $k$ ثابت ($k\ll n$):
\[
\binom nk p^k(1-p)^{n-k}\ \longrightarrow\ \frac{\lambda^k}{k!}e^{-\lambda}.
\]
\end{theorem}
\begin{proof}
\[
\binom nk p^k(1-p)^{n-k}=\frac{n(n-1)\cdots(n-k+1)}{k!}\,p^k(1-p)^{n-k}.
\]
سه عامل را جداگانه بررسی می‌کنیم:
\begin{itemize}
\item $\dfrac{n(n-1)\cdots(n-k+1)}{n^k}=1\cdot\big(1-\tfrac1n\big)\cdots\big(1-\tfrac{k-1}{n}\big)\to1$؛ پس صورت تقریباً $n^k$ است و $n^kp^k=(np)^k=\lambda^k$.
\item $(1-p)^{-k}\to1$ چون $p\to0$ و $k$ ثابت است؛ پس $(1-p)^{n-k}\simeq(1-p)^n$.
\item $(1-p)^n=\big(1-\frac\lambda n\big)^n\to e^{-\lambda}$ (حد معروف).
\end{itemize}
پس
\[
\simeq\frac{n^kp^k}{k!}(1-p)^n\simeq\frac{(np)^k}{k!}e^{-np},
\]
که در حد، دقیقاً برابر با این مقدار خواهد شد.
\end{proof}

\begin{note}
\notetitle \textbf{تعبیر:} پواسون مدل «تعداد زیادی فرصت مستقل، هر یک با احتمال بسیار کم» است: مثلاً یک ساعت را به $n$ بازه بسیار کوچک تقسیم کنید که در هر کدام با احتمال $p=\lambda/n$ یک تماس رخ می‌دهد. وقتی $n\to\infty$، تعداد تماس‌ها پواسون با میانگین $\lambda$ می‌شود. به همین دلیل پواسون برای رخدادهای «نادر» و «بی‌حافظه» در زمان به کار می‌رود.
\end{note}

\begin{example}[دقت تقریب پواسون]
در یک خط تولید هر قطعه با احتمال $0.002$ معیوب است. در یک محموله ۱۰۰۰تایی، احتمال اینکه دقیقاً ۳ قطعه معیوب باشد را به صورت دقیق و با تقریب پواسون حساب کنید.
\begin{solution}
مقدار دقیق (دوجمله‌ای): $\binom{1000}{3}(0.002)^3(0.998)^{997}\approx0.1806$.
تقریب پواسون با $\lambda=np=2$: $e^{-2}\frac{2^3}{3!}\approx0.1804$.
جدول زیر نشان می‌دهد تقریب برای همه $k$ها بسیار دقیق است:
\begin{center}
\renewcommand{\arraystretch}{1.2}
\begin{tabular}{c|cccccc}
$k$ & $0$ & $1$ & $2$ & $3$ & $4$ & $5$\\\hline
دوجمله‌ای & $0.1351$ & $0.2707$ & $0.2709$ & $0.1806$ & $0.0902$ & $0.0360$\\
پواسون & $0.1353$ & $0.2707$ & $0.2707$ & $0.1804$ & $0.0902$ & $0.0361$
\end{tabular}
\end{center}
\end{solution}
\end{example}

\subsection{توزیع هندسی \lr{(Geometric)} ــ تکمیلی}
اگر آزمون‌های برنولی مستقل با احتمال موفقیت $p$ را تا رسیدن به \emph{اولین} موفقیت تکرار کنیم، تعداد آزمون‌ها $X$ توزیع هندسی دارد:
\[
P_X(k)=(1-p)^{k-1}p,\qquad k=1,2,\ldots
\]
زیرا باید $k-1$ شکست متوالی و سپس یک موفقیت رخ دهد. جمع احتمال‌ها یک سری هندسی است: $p\sum_{k\ge1}(1-p)^{k-1}=p\cdot\frac1{p}=1$. این توزیع همتای گسسته توزیع نمایی (جلسه بعد) و تنها توزیع گسسته \textbf{بی‌حافظه} است.

\section{تمرین‌های تکمیلی}

\begin{exercise}
ثابت $c$ را چنان بیابید که $P_X(k)=c\,(1/3)^k$ برای $k=0,1,2,\ldots$ یک \lr{PMF} باشد. سپس $P(X\ge2)$ را حساب کنید.
\end{exercise}
\begin{solution}
$\sum_{k\ge0}(1/3)^k=\frac{1}{1-1/3}=\frac32$؛ پس $c=\frac23$. آنگاه
$P(X\ge2)=1-P(0)-P(1)=1-\frac23-\frac29=\frac19$.
\end{solution}

\begin{exercise}
به یک سرور به طور متوسط ۳ درخواست در دقیقه می‌رسد (مدل پواسون). احتمال اینکه در یک دقیقه مشخص هیچ درخواستی نرسد، و احتمال اینکه بیش از ۲ درخواست برسد چقدر است؟
\end{exercise}
\begin{solution}
$X\sim\mathrm{Pois}(3)$. $P(X=0)=e^{-3}\approx0.0498$.
$P(X>2)=1-e^{-3}\big(1+3+\frac92\big)=1-8.5e^{-3}\approx1-0.4232=0.5768$.
\end{solution}

\begin{exercise}
یک بسته داده از طریق کانالی ارسال می‌شود که هر بیت را مستقلاً با احتمال $10^{-3}$ خراب می‌کند. بسته ۲۰۰۰ بیتی است و کد تصحیح خطا تا ۲ بیت خطا را تصحیح می‌کند. احتمال دریافت صحیح بسته را تقریب بزنید.
\end{exercise}
\begin{solution}
تعداد بیت‌های خراب $X\sim\mathrm{Binomial}(2000,10^{-3})\approx\mathrm{Pois}(2)$. احتمال موفقیت:
$P(X\le2)\approx e^{-2}(1+2+2)=5e^{-2}\approx0.677$.
\end{solution}

\begin{exercise}
تاسی را آن‌قدر می‌اندازیم تا ۶ بیاید. احتمال اینکه بیش از ۱۰ پرتاب لازم باشد چقدر است؟
\end{exercise}
\begin{solution}
تعداد پرتاب‌ها هندسی با $p=\frac16$ است. «بیش از ۱۰ پرتاب لازم باشد» یعنی ۱۰ پرتاب اول همگی شکست باشند: $P(X>10)=(5/6)^{10}\approx0.1615$.
\end{solution}
